% Unit 5 (Kinetics) guided notes -- 7 blocks, CED sequence.
\documentclass{apchem-notes}

\apcunit{5}
\apcunittitle{Kinetics}
\apcdoctitle{Guided Notes}
\apcsubtitle{CED 5.1--5.11 \textbullet\ Zumdahl Ch 12 \textbullet\ 7 blocks}

\begin{document}
\apctitleblock

\begin{apcnote}[The one idea this unit runs on]
Thermodynamics (Unit 9) will tell you \emph{whether} a reaction is favored.
Equilibrium (Unit 7) will tell you \emph{how far} it goes. Neither says
anything about \emph{how fast}.

Rate is a separate question with a separate answer, and the answer is
always experimental. \textbf{You cannot read a rate law off a balanced
equation.} Zumdahl's own example: $\ce{2N2O5 -> 4NO2 + O2}$ has a
coefficient of 2 and is \emph{first} order in \ce{N2O5}.

The single exception is an \term{elementary step} --- one actual molecular
event --- where the molecularity does give the order. Recognizing which case
you are in is the skill this unit is really testing.
\end{apcnote}

% =====================================================================
\blocklesson{Reaction Rates}{CED 5.1 \textbullet\ Zumdahl \S12.1}

\begin{objectives}
  \item Express a rate in terms of any species in the equation.
  \item Distinguish average, instantaneous, and initial rate.
\end{objectives}

\reading{Zumdahl \S12.1}{594--598}

\begin{warmup}[4]
  \item Units of concentration: \answerblank[2.6cm]{mol/L (M)}
  \item Slope of a graph is: \answerblank[3.4cm]{rise over run}
  \item In \ce{2A -> B}, A is consumed how fast relative to B?
        \answerblank[3cm]{twice as fast}
  \item Rate must always be a \answerblank[2.6cm]{positive} quantity.
\end{warmup}

\segment{INSTRUCTION A}{25}{Defining a rate}

\notesection{Rate is a change per unit time}{5.1}
\practice{4}

\term{Reaction rate} is the change in concentration of a species per unit
time. Because reactants disappear, a minus sign is used to keep the rate
\answerblank[2cm]{positive}:
\[ \text{rate} = -\frac{\Delta[\text{reactant}]}{\Delta t}
   = +\frac{\Delta[\text{product}]}{\Delta t} \]

\begin{apctrap}
\textbf{Divide by the coefficient.} For $\ce{aA + bB -> cC + dD}$,
\[ \text{rate} = -\frac{1}{a}\frac{\Delta[\ce{A}]}{\Delta t}
   = -\frac{1}{b}\frac{\Delta[\ce{B}]}{\Delta t}
   = +\frac{1}{c}\frac{\Delta[\ce{C}]}{\Delta t}
   = +\frac{1}{d}\frac{\Delta[\ce{D}]}{\Delta t} \]
Without the coefficients, ``the rate'' has a different value depending on
which species you watched --- which is exactly the confusion the convention
exists to prevent. This is the most common lost point in Block 1.
\end{apctrap}

\begin{workedexample}[Worked example 1: one reaction, three rates]
For $\ce{2N2O5(g) -> 4NO2(g) + O2(g)}$, oxygen is forming at
\SI{2.0e-3}{\Molar\per\second}. Find the rate of the reaction and the rate
of disappearance of \ce{N2O5}.

Oxygen has coefficient 1, so the \emph{reaction rate} is
$\SI{2.0e-3}{\Molar\per\second}$.

\ce{N2O5} has coefficient 2, so it disappears \emph{twice} as fast:
\[ -\frac{\Delta[\ce{N2O5}]}{\Delta t} = 2(2.0\times10^{-3})
   = \mathbf{4.0\times10^{-3}}~\si{\Molar\per\second} \]
and \ce{NO2} appears four times as fast,
$\SI{8.0e-3}{\Molar\per\second}$.
\end{workedexample}

\segment{GUIDED PRACTICE}{15}{Three kinds of rate}

A plot of $[\text{reactant}]$ against time is a curve that starts steep and
flattens.

\begin{enumerate}[leftmargin=*,itemsep=0.35em]
  \item \term{Average rate} over an interval is the slope of the
        \answerblank[3.4cm]{straight line} between two points.
  \item \term{Instantaneous rate} at one moment is the slope of the
        \answerblank[2.4cm]{tangent} there.
  \item \term{Initial rate} is the instantaneous rate at
        \answerblank[1.8cm]{$t = 0$}.
  \item Why does the curve flatten? \answerlines[2]{As reactant is consumed
        its concentration falls, so collisions between reactant particles
        become less frequent and the rate decreases.}
  \item Why do rate-law experiments use the \emph{initial} rate?
        \answerlines[2]{At $t = 0$ no products have accumulated, so no
        reverse reaction and no product interference complicate the
        measurement.}
\end{enumerate}

\segment{APPLICATION}{20}{Rate conversions}

\begin{enumerate}[leftmargin=*,itemsep=0.4em]
  \item For $\ce{N2 + 3H2 -> 2NH3}$, \ce{H2} is consumed at
        \SI{0.030}{\Molar\per\second}. Find the rate of the reaction and the
        rate of formation of \ce{NH3}. \workspace[2.4cm]
        \keyonly{\begin{apcanswer}rate $= (1/3)(0.030) =
        \mathbf{0.010}~\si{\Molar\per\second}$; \ce{NH3} forms at
        $2(0.010) = \mathbf{0.020}~\si{\Molar\per\second}$\end{apcanswer}}
  \item A student reports ``the rate is
        \SI{0.030}{\Molar\per\second}'' after watching \ce{H2}. What is
        missing from that statement? \answerlines[2]{Which species was
        measured. Without dividing by the coefficient, the number is a rate
        of disappearance of \ce{H2}, not the reaction rate --- those differ
        by a factor of 3 here.}
\end{enumerate}

\begin{exitticket}
For $\ce{2A -> 3B}$, if B appears at \SI{6.0e-4}{\Molar\per\second}, how
fast does A disappear?
\keyonly{\begin{apcanswer}Reaction rate $= (1/3)(6.0\times10^{-4}) =
2.0\times10^{-4}$. A disappears at $2(2.0\times10^{-4}) =
\mathbf{4.0\times10^{-4}}~\si{\Molar\per\second}$.\end{apcanswer}}
\end{exitticket}

% =====================================================================
\blocklesson{Rate Laws and Orders}{CED 5.2 \textbullet\ Zumdahl \S12.2--12.3}

\begin{objectives}
  \item Write a rate law and identify reaction orders.
  \item Determine orders and $k$ from initial-rate data.
\end{objectives}

\reading{Zumdahl \S12.2--12.3}{598--604}

\begin{warmup}[3]
  \item Rate for $\ce{2A -> B}$ in terms of A:
        \answerblank[4cm]{$-\tfrac{1}{2}\Delta[\ce{A}]/\Delta t$}
  \item Initial rate is measured at: \answerblank[2cm]{$t = 0$}
  \item Orders come from: \answerblank[2.6cm]{experiment}
\end{warmup}

\segment{INSTRUCTION A}{25}{The form of a rate law}

\notesection{Rate law, orders, and the rate constant}{5.2}
\practice{5}

\[ \text{rate} = k[\ce{A}]^m[\ce{B}]^n \]

\begin{itemize}[leftmargin=*,itemsep=0.3em]
  \item $k$ is the \term{rate constant}: constant with concentration,
        \emph{not} with \answerblank[2.7cm]{temperature}.
  \item $m$ and $n$ are the \term{orders} --- determined
        \answerblank[3cm]{experimentally}, never from coefficients.
  \item The \term{overall order} is \answerblank[2cm]{$m + n$}.
\end{itemize}

\begin{apctrap}
\textbf{Orders are not coefficients.} This is the single most-tested
misconception in the unit. Zumdahl's counterexample is worth memorizing:
$\ce{2N2O5 -> 4NO2 + O2}$ carries a coefficient of 2 and is experimentally
\emph{first} order in \ce{N2O5}.

A related trap: the units of $k$ change with the overall order, so they are
a free check on your work. First order gives
\si{\per\second}; second order \si{\per\Molar\per\second}; zero order
\si{\Molar\per\second}.
\end{apctrap}

\notesection{The method of initial rates}{5.2}

Change one concentration at a time and see what the rate does. If doubling
$[\ce{A}]$ multiplies the rate by $2^m$, then $m$ is the order in A.

\begin{center}
\begin{tabular}{@{}lll@{}}
\hline
\textbf{Double [A] and the rate\dots} & \textbf{order in A} & \\
\hline
does not change & 0 & $2^0 = 1$ \\
doubles & 1 & $2^1 = 2$ \\
quadruples & 2 & $2^2 = 4$ \\
\hline
\end{tabular}
\end{center}

\begin{workedexample}[Worked example 2: reading orders from a table]
\begin{center}
\begin{tabular}{@{}cccc@{}}
\hline
\textbf{Exp} & $[\ce{A}]_0$ & $[\ce{B}]_0$ &
  \textbf{initial rate} (\si{\Molar\per\second}) \\
\hline
1 & 0.10 & 0.10 & $2.0\times10^{-3}$ \\
2 & 0.20 & 0.10 & $8.0\times10^{-3}$ \\
3 & 0.10 & 0.20 & $4.0\times10^{-3}$ \\
\hline
\end{tabular}
\end{center}

\textbf{Order in A} --- compare 1 and 2, where only $[\ce{A}]$ changed.
Doubling A multiplied the rate by 4, so $m = \mathbf{2}$.

\textbf{Order in B} --- compare 1 and 3. Doubling B doubled the rate, so
$n = \mathbf{1}$.

\[ \text{rate} = k[\ce{A}]^2[\ce{B}] \qquad\text{overall order } 3 \]

\textbf{Find $k$} from any row:
\[ k = \frac{2.0\times10^{-3}}{(0.10)^2(0.10)}
     = \mathbf{2.0}~\si{\per\Molar\squared\per\second} \]
Units of $\si{\Molar^{-2}\second^{-1}}$ confirm third order overall.
\end{workedexample}

\segment{GUIDED PRACTICE}{15}{A second data set}

\begin{center}
\begin{tabular}{@{}cccc@{}}
\hline
\textbf{Exp} & $[\ce{X}]_0$ & $[\ce{Y}]_0$ &
  \textbf{rate} (\si{\Molar\per\second}) \\
\hline
1 & 0.20 & 0.10 & $1.5\times10^{-2}$ \\
2 & 0.40 & 0.10 & $3.0\times10^{-2}$ \\
3 & 0.20 & 0.20 & $1.5\times10^{-2}$ \\
\hline
\end{tabular}
\end{center}

\begin{enumerate}[leftmargin=*,itemsep=0.35em]
  \item Order in X: \answerblank[1.6cm]{1}
  \item Order in Y: \answerblank[1.6cm]{0}
  \item Rate law: \answerblank[3.4cm]{rate $= k[\ce{X}]$}
  \item $k = $ \answerblank[3cm]{\SI{0.075}{\per\second}}
  \item What does zero order in Y mean physically?
        \answerlines[2]{Changing $[\ce{Y}]$ has no effect on the rate ---
        Y does not appear in the rate-determining step, even though it is a
        reactant in the overall equation.}
\end{enumerate}

\segment{APPLICATION}{20}{Using a rate law}

\begin{enumerate}[leftmargin=*,itemsep=0.4em]
  \item For rate $= k[\ce{A}]^2[\ce{B}]$ with
        $k = \SI{2.0}{\per\Molar\squared\per\second}$, find the rate when
        $[\ce{A}] = 0.30$ and $[\ce{B}] = 0.20$. \workspace[2cm]
        \keyonly{\begin{apcanswer}$(2.0)(0.30)^2(0.20) = (2.0)(0.090)(0.20)
        = \mathbf{0.036}~\si{\Molar\per\second}$\end{apcanswer}}
  \item In that same reaction, what happens to the rate if $[\ce{A}]$ is
        tripled and $[\ce{B}]$ halved? \answerlines[2]{Rate changes by
        $3^2 \times \tfrac{1}{2} = 4.5$ --- it increases by a factor of
        4.5.}
  \item A student writes rate $= k[\ce{A}]^2[\ce{B}]$ for
        $\ce{2A + B -> C}$ ``because of the coefficients.'' Explain why the
        reasoning is wrong even though the answer happens to match.
        \answerlines[3]{Orders must be measured, not inferred. Coefficients
        describe overall stoichiometry; orders describe the mechanism's
        slow step. They coincide only by chance unless the reaction happens
        to be a single elementary step.}
\end{enumerate}

\begin{exitticket}
Give the units of $k$ for a zero-order, first-order, and second-order
reaction, and say why the units are a useful check.
\keyonly{\begin{apcanswer}Zero: \si{\Molar\per\second}. First:
\si{\per\second}. Second: \si{\per\Molar\per\second}. Because rate always
has units of \si{\Molar\per\second}, the units of $k$ are fixed by the
overall order --- so if your $k$ has the wrong units, your orders are
wrong.\end{apcanswer}}
\end{exitticket}

% =====================================================================
\blocklesson{Concentration Over Time}{CED 5.3 \textbullet\ Zumdahl \S12.4}

\begin{objectives}
  \item Identify order from which plot is linear.
  \item Apply integrated rate laws and half-life.
\end{objectives}

\reading{Zumdahl \S12.4}{604--615}

\begin{warmup}[3]
  \item A rate law gives rate in terms of:
        \answerblank[3cm]{concentration}
  \item Order in Y is 0 means the rate:
        \answerblank[4.6cm]{does not depend on Y}
  \item Units of $k$ for first order:
        \answerblank[2.4cm]{\si{\per\second}}
\end{warmup}

\segment{INSTRUCTION A}{25}{Three orders, three straight lines}

\notesection{Integrated rate laws}{5.3}
\practice{5}

A rate law tells you the rate \emph{now}; the integrated form tells you the
concentration \emph{at time $t$}. Each order gives a different straight-line
plot, and that is how order is identified from data.

\begin{center}
\renewcommand{\arraystretch}{1.3}
\begin{tabular}{@{}llll@{}}
\hline
\textbf{Order} & \textbf{Integrated law} & \textbf{Linear plot} &
  \textbf{Half-life} \\
\hline
0 & $[\ce{A}] = -kt + [\ce{A}]_0$ & $[\ce{A}]$ vs $t$ &
  $t_{1/2} = [\ce{A}]_0/2k$ \\
1 & $\ln[\ce{A}] = -kt + \ln[\ce{A}]_0$ & $\ln[\ce{A}]$ vs $t$ &
  $t_{1/2} = 0.693/k$ \\
2 & $1/[\ce{A}] = kt + 1/[\ce{A}]_0$ & $1/[\ce{A}]$ vs $t$ &
  $t_{1/2} = 1/(k[\ce{A}]_0)$ \\
\hline
\end{tabular}
\end{center}

In every case the slope is $\pm k$ --- negative for zero and first order,
positive for \answerblank[2.6cm]{second order}.

\begin{apctrap}
\textbf{``The graph is linear'' is not an answer.} You must say
\emph{which quantity} was plotted. ``The plot of $\ln[\ce{A}]$ against time
is linear, therefore the reaction is first order'' earns the point;
``the graph is a straight line'' does not.

\textbf{Only first-order half-life is constant.} $t_{1/2} = 0.693/k$
contains no $[\ce{A}]_0$, so a first-order reaction takes the same time to
halve no matter where you start. For zero and second order the half-life
depends on concentration --- so a constant half-life is itself evidence of
first order.
\end{apctrap}

\begin{workedexample}[Worked example 3: a first-order decay]
$k = \SI{0.0250}{\per\second}$ and $[\ce{A}]_0 = \SI{0.500}{\Molar}$.

\textbf{Concentration after \SI{60.0}{\second}:}
\[ \ln[\ce{A}] = -(0.0250)(60.0) + \ln(0.500) = -1.50 - 0.693 = -2.193 \]
\[ [\ce{A}] = e^{-2.193} = \mathbf{0.112}~\si{\Molar} \]

\textbf{Half-life:} $t_{1/2} = 0.693/0.0250 = \mathbf{27.7}~\si{\second}$.

\textbf{Sanity check:} two half-lives is \SI{55.4}{\second}, at which point
$[\ce{A}]$ should be $0.125$~M. Our \SI{60}{\second} value of $0.112$ is a
little below that --- consistent.
\end{workedexample}

\segment{GUIDED PRACTICE}{15}{Second and zero order}

\begin{enumerate}[leftmargin=*,itemsep=0.4em]
  \item Second order, $k = \SI{0.15}{\per\Molar\per\second}$,
        $[\ce{A}]_0 = \SI{0.20}{\Molar}$. Find $[\ce{A}]$ after
        \SI{100}{\second}. \workspace[2.2cm]
        \keyonly{\begin{apcanswer}$1/[\ce{A}] = (0.15)(100) + 1/0.20 =
        15 + 5.0 = 20$; $[\ce{A}] =
        \mathbf{0.050}~\si{\Molar}$\end{apcanswer}}
  \item Its half-life: \answerblank[2.6cm]{\SI{33}{\second}}
  \item Zero order, $k = \SI{0.020}{\Molar\per\second}$,
        $[\ce{A}]_0 = \SI{1.00}{\Molar}$. Find $t_{1/2}$ and
        $[\ce{A}]$ after \SI{30}{\second}. \workspace[2cm]
        \keyonly{\begin{apcanswer}$t_{1/2} = 1.00/(2 \times 0.020) =
        \mathbf{25}~\si{\second}$;
        $[\ce{A}] = 1.00 - (0.020)(30) =
        \mathbf{0.40}~\si{\Molar}$\end{apcanswer}}
\end{enumerate}

\segment{APPLICATION}{20}{Identifying order from data}

\begin{enumerate}[leftmargin=*,itemsep=0.4em]
  \item A reaction's half-life is \SI{25.0}{\minute} at the start and
        \SI{25.0}{\minute} again after half the reactant is gone. What is
        the order, and what is $k$? \workspace[2cm]
        \keyonly{\begin{apcanswer}Constant half-life $\Rightarrow$
        \textbf{first order}. $k = 0.693/25.0 =
        \mathbf{0.0277}~\si{\per\minute}$\end{apcanswer}}
  \item Data for a decomposition give a curved plot of $[\ce{A}]$ vs $t$, a
        curved plot of $\ln[\ce{A}]$ vs $t$, and a straight plot of
        $1/[\ce{A}]$ vs $t$ with slope $+0.42$. State the order, $k$, and
        the rate law. \answerlines[3]{Second order, since only
        $1/[\ce{A}]$ vs $t$ is linear. The slope equals $k$, so
        $k = \SI{0.42}{\per\Molar\per\second}$ and
        rate $= k[\ce{A}]^2$.}
  \item Why is a constant half-life such strong evidence of first order?
        \answerlines[2]{Because $t_{1/2} = 0.693/k$ is the only one of the
        three expressions that does not contain $[\ce{A}]_0$. For zero and
        second order the half-life changes as the reaction proceeds.}
\end{enumerate}

\begin{exitticket}
Which plot is linear for each order, and what is the slope in each case?
\keyonly{\begin{apcanswer}Zero order: $[\ce{A}]$ vs $t$, slope $-k$. First
order: $\ln[\ce{A}]$ vs $t$, slope $-k$. Second order: $1/[\ce{A}]$ vs $t$,
slope $+k$.\end{apcanswer}}
\end{exitticket}

% =====================================================================
\blocklesson{Elementary Steps and Mechanisms}{CED 5.4, 5.7, 5.8 \textbullet\ Zumdahl \S12.5}

\begin{objectives}
  \item Write the rate law for an elementary step from its molecularity.
  \item Test whether a proposed mechanism is acceptable.
\end{objectives}

\reading{Zumdahl \S12.5}{615--620}

\begin{warmup}[3]
  \item Orders normally come from:
        \answerblank[2.6cm]{experiment}
  \item First-order half-life $=$ \answerblank[2.4cm]{$0.693/k$}
  \item Second order gives which linear plot?
        \answerblank[3cm]{$1/[\ce{A}]$ vs $t$}
\end{warmup}

\segment{INSTRUCTION A}{25}{The one place coefficients do give orders}

\notesection{Elementary steps}{5.4}
\practice{5}

A \term{mechanism} is the sequence of \term{elementary steps} --- single
molecular events --- that add up to the overall reaction.

\begin{center}
\fbox{\parbox{0.88\linewidth}{\centering
For an \textbf{elementary step only}, the coefficients \emph{are} the
orders. \\
This is the only exception in the entire unit.}}
\end{center}

\term{Molecularity} is the number of particles colliding in that step:

\begin{center}
\begin{tabular}{@{}lll@{}}
\hline
\textbf{Molecularity} & \textbf{Step} & \textbf{Rate law} \\
\hline
unimolecular & $\ce{A -> products}$ & rate $= k[\ce{A}]$ \\
bimolecular & $\ce{A + B -> products}$ & rate $= k[\ce{A}][\ce{B}]$ \\
bimolecular & $\ce{2A -> products}$ & rate $= k[\ce{A}]^2$ \\
termolecular & $\ce{A + B + C -> products}$ & rate $= k[\ce{A}][\ce{B}][\ce{C}]$ \\
\hline
\end{tabular}
\end{center}

Termolecular steps are \answerblank[2.4cm]{very rare} --- three particles
colliding simultaneously with the right orientation is improbable.

\notesection{Two species with special roles}{5.7}

\begin{itemize}[leftmargin=*,itemsep=0.3em]
  \item An \term{intermediate} is
        \answerblank[5.2cm]{produced then consumed} --- it does not appear
        in the overall equation, and it must \textbf{never} appear in a
        final rate law.
  \item A \term{catalyst} is \answerblank[5.6cm]{consumed then regenerated}
        --- it appears in the first step and is returned in a later one.
\end{itemize}

\segment{INSTRUCTION B}{20}{Testing a mechanism}

\notesection{The two requirements}{5.8}
\practice{6}

A proposed mechanism is acceptable only if \textbf{both} hold:

\begin{enumerate}[leftmargin=*,itemsep=0.3em]
  \item The steps sum to the \answerblank[5.2cm]{overall balanced equation}
        (intermediates cancel).
  \item Its predicted rate law matches the
        \answerblank[4.4cm]{experimental rate law}.
\end{enumerate}

Even then the mechanism is only \emph{consistent} with the data --- never
proven. Several mechanisms can predict the same rate law.

\begin{workedexample}[Worked example 4: slow first step]
Overall: $\ce{2NO2 + F2 -> 2NO2F}$, experimentally
rate $= k[\ce{NO2}][\ce{F2}]$.

Proposed:
\begin{align*}
  \text{Step 1 (slow):} &\quad \ce{NO2 + F2 -> NO2F + F} \\
  \text{Step 2 (fast):} &\quad \ce{NO2 + F -> NO2F}
\end{align*}

\textbf{Do they sum?} Adding gives
$\ce{2NO2 + F2 -> 2NO2F}$, with the F atom cancelling. \checkmark

\textbf{Predicted rate law?} The slow step is
\term{rate-determining}, and it is elementary, so its rate law is
rate $= k[\ce{NO2}][\ce{F2}]$. \checkmark\ Matches.

Note that F is an \textbf{intermediate} --- made in step 1, used in step 2 ---
and correctly does not appear in the rate law.
\end{workedexample}

\begin{apctrap}
\textbf{When the slow step is first, the rate law is just that step's.}
That is the easy case, and it is the one most exam questions use. The hard
case --- a \emph{fast} first step --- is Block 6.

If your derived rate law contains an intermediate, you are not finished.
\end{apctrap}

\segment{APPLICATION}{20}{Evaluating mechanisms}

\begin{enumerate}[leftmargin=*,itemsep=0.4em]
  \item Overall $\ce{2NO + O2 -> 2NO2}$ has experimental
        rate $= k[\ce{NO}]^2[\ce{O2}]$. Is a single termolecular step a
        plausible mechanism? \answerlines[3]{It would predict the correct
        rate law, but a simultaneous three-body collision is extremely
        improbable, so a two-step mechanism with the same overall rate law
        is far more plausible. Matching the rate law makes a mechanism
        consistent, not correct.}
  \item For the mechanism $\ce{A2 -> 2A}$ (slow),
        $\ce{A + B -> AB}$ (fast), give the overall equation and predicted
        rate law. \answerlines[3]{Doubling the second step to cancel the
        intermediate: overall $\ce{A2 + 2B -> 2AB}$. The slow step is
        elementary and unimolecular, so rate $= k[\ce{A2}]$ --- note that B
        does not appear at all.}
  \item Identify the intermediate and the catalyst:
        $\ce{X + C -> XC}$, then $\ce{XC + Y -> XY + C}$.
        \answerblank[6cm]{XC intermediate; C catalyst}
\end{enumerate}

\begin{exitticket}
State the two tests a proposed mechanism must pass, and say why passing
both still does not prove it.
\keyonly{\begin{apcanswer}The steps must sum to the overall balanced
equation, and the predicted rate law must match the experimental one.
Passing both makes the mechanism \emph{consistent} with the evidence, but
other mechanisms can predict the same rate law, so it is never
proven.\end{apcanswer}}
\end{exitticket}

% =====================================================================
\blocklesson{Collision Model and Energy Profiles}{CED 5.5, 5.6 \textbullet\ Zumdahl \S12.6}

\begin{objectives}
  \item Explain rate in terms of collision frequency, energy, and
        orientation.
  \item Read and draw a reaction energy profile.
\end{objectives}

\reading{Zumdahl \S12.6}{620--626}

\begin{warmup}[3]
  \item Molecularity gives the order only for:
        \answerblank[3.8cm]{an elementary step}
  \item An intermediate is: \answerblank[5.2cm]{produced then consumed}
  \item Raising temperature does what to rate?
        \answerblank[2.6cm]{increases it}
\end{warmup}

\segment{INSTRUCTION A}{25}{Why most collisions do nothing}

\notesection{Three requirements for a reactive collision}{5.5}
\practice{6}

\begin{enumerate}[leftmargin=*,itemsep=0.3em]
  \item The particles must \answerblank[2.4cm]{collide}.
  \item They must collide with at least the
        \answerblank[3.6cm]{activation energy}, $E_a$.
  \item They must collide with the correct
        \answerblank[2.6cm]{orientation}.
\end{enumerate}

Only a tiny fraction of collisions satisfy all three, which is why rates are
so much smaller than collision frequencies.

\notesection{Temperature and the Arrhenius equation}{5.5}

\[ k = Ae^{-E_a/RT} \]

Raising the temperature does \emph{not} mainly work by making collisions
more frequent --- that effect is small. It works because the distribution of
molecular kinetic energies \answerblank[2.4cm]{broadens}, so a much larger
fraction of collisions clears $E_a$.

\begin{workedexample}[Worked example 5: a ten-degree rise]
A reaction has $E_a = \SI{50.0}{\kilo\joule\per\mole}$. By what factor does
$k$ change from \SI{300}{\kelvin} to \SI{310}{\kelvin}?

\[ \ln\frac{k_2}{k_1} = -\frac{E_a}{R}\left(\frac{1}{T_2} -
   \frac{1}{T_1}\right)
   = -\frac{50000}{8.314}\left(\frac{1}{310} - \frac{1}{300}\right)
   = 0.647 \]
\[ \frac{k_2}{k_1} = e^{0.647} = \mathbf{1.9} \]

A \SI{10}{\kelvin} rise nearly \textbf{doubles} the rate --- the familiar
rule of thumb, and now derived rather than asserted.

With $E_a = \SI{100}{\kilo\joule\per\mole}$ the same rise multiplies $k$ by
\textbf{3.6}: the larger the barrier, the more temperature-sensitive the
reaction.
\end{workedexample}

\segment{INSTRUCTION B}{20}{Reading the energy profile}

\notesection{What the diagram shows}{5.6}
\practice{3}

\begin{center}
\begin{tikzpicture}
\begin{axis}[width=10cm, height=5.2cm,
    xlabel={reaction progress}, ylabel={potential energy},
    xtick=\empty, ytick=\empty, xmin=0, xmax=10, ymin=0, ymax=10,
    axis lines=left]
  \addplot[seriesA, samples=200, domain=0:10]
    {2 + 6*exp(-((x-5)^2)/2.2)};
  \draw[densely dashed] (axis cs:0,2) -- (axis cs:5,2);
  \draw[densely dashed] (axis cs:5,8) -- (axis cs:10,8);
  \draw[densely dashed] (axis cs:10,3.6) -- (axis cs:5,3.6);
  \draw[<->,thick] (axis cs:1.6,2) -- (axis cs:1.6,8)
    node[midway,right,font=\scriptsize] {$E_a$ (forward)};
  \draw[<->,thick] (axis cs:8.6,3.6) -- (axis cs:8.6,8)
    node[midway,right,font=\scriptsize] {$E_a$ (reverse)};
  \draw[<->,thick] (axis cs:0.5,2) -- (axis cs:0.5,3.6)
    node[midway,left,font=\scriptsize] {$\Delta H$};
  \node[font=\scriptsize,anchor=north] at (axis cs:0.9,1.9) {reactants};
  \node[font=\scriptsize,anchor=north] at (axis cs:9.2,3.5) {products};
  \node[font=\scriptsize,anchor=south] at (axis cs:5,8.1)
    {transition state};
\end{axis}
\end{tikzpicture}
\end{center}

\begin{itemize}[leftmargin=*,itemsep=0.3em]
  \item $E_a$ is measured from the \emph{reactants} up to the
        \answerblank[3.4cm]{transition state}.
  \item $\Delta H$ is the difference between
        \answerblank[4.6cm]{products and reactants} --- independent of the
        barrier height.
  \item $E_a(\text{forward}) - E_a(\text{reverse}) = $
        \answerblank[1.8cm]{$\Delta H$}.
\end{itemize}

\begin{apctrap}
$E_a$ and $\Delta H$ are independent. A strongly exothermic reaction can
have an enormous barrier and be immeasurably slow --- diamond to graphite is
the standard example. Do not use ``very exothermic'' as evidence that a
reaction is fast, or a large $E_a$ as evidence that it is endothermic.
\end{apctrap}

\segment{APPLICATION}{20}{Profiles and barriers}

\begin{enumerate}[leftmargin=*,itemsep=0.4em]
  \item A reaction has $E_a(\text{fwd}) = \SI{80}{\kilo\joule}$ and
        $E_a(\text{rev}) = \SI{120}{\kilo\joule}$. Find $\Delta H$ and
        classify the reaction. \answerblank[5cm]{$-40$ kJ; exothermic}
  \item Another has $E_a(\text{fwd}) = \SI{150}{\kilo\joule}$ and
        $\Delta H = \SI{+60}{\kilo\joule}$. Find $E_a(\text{rev})$.
        \answerblank[3cm]{\SI{90}{\kilo\joule}}
  \item Explain, in terms of the energy distribution, why raising the
        temperature by \SI{10}{\kelvin} can nearly double a rate when it
        raises the average kinetic energy by only about 3\%.
        \answerlines[3]{Rate depends on the \emph{fraction} of collisions
        exceeding $E_a$, which lies far out in the tail of the distribution.
        Broadening the distribution slightly moves a disproportionately
        large number of particles past that threshold, because the tail is
        where the change is proportionally greatest.}
\end{enumerate}

\begin{exitticket}
Two reactions have the same $\Delta H$ but different $E_a$. What differs
between them, and what does not?
\keyonly{\begin{apcanswer}The rates differ --- the one with the smaller
$E_a$ is faster, because a larger fraction of collisions clears its barrier.
The thermodynamics is identical: same $\Delta H$, same relative energies of
reactants and products, and (at a given temperature) the same
$K$.\end{apcanswer}}
\end{exitticket}

% =====================================================================
\blocklesson{Fast First Steps and Multistep Profiles}{CED 5.9, 5.10 \textbullet\ Zumdahl \S12.5--12.6}

\begin{objectives}
  \item Derive a rate law when the first step is a fast equilibrium.
  \item Interpret a multistep energy profile.
\end{objectives}

\reading{Zumdahl \S12.5--12.6}{615--626}

\begin{warmup}[3]
  \item A final rate law must never contain:
        \answerblank[3.2cm]{an intermediate}
  \item $E_a(\text{fwd}) - E_a(\text{rev}) = $
        \answerblank[1.8cm]{$\Delta H$}
  \item The slow step is called the:
        \answerblank[4.4cm]{rate-determining step}
\end{warmup}

\segment{INSTRUCTION A}{25}{When the slow step is not first}

\notesection{The pre-equilibrium approximation}{5.9}
\practice{5}

If the \emph{first} step is fast and reversible and the second is slow, the
first step reaches equilibrium before the second consumes much of the
intermediate. Write the rate law for the slow step, then use the
equilibrium of step 1 to \answerblank[5.4cm]{eliminate the intermediate}.

\begin{workedexample}[Worked example 6: eliminating an intermediate]
Overall $\ce{2NO + O2 -> 2NO2}$, experimental
rate $= k[\ce{NO}]^2[\ce{O2}]$.

\begin{align*}
  \text{Step 1 (fast, equilibrium):} &\quad \ce{2NO <=> N2O2} \\
  \text{Step 2 (slow):} &\quad \ce{N2O2 + O2 -> 2NO2}
\end{align*}

\textbf{Start with the slow step:}
rate $= k_2[\ce{N2O2}][\ce{O2}]$. But \ce{N2O2} is an intermediate, so this
is not yet a valid rate law.

\textbf{Use the equilibrium:} for step 1,
$K_1 = \dfrac{[\ce{N2O2}]}{[\ce{NO}]^2}$, so
$[\ce{N2O2}] = K_1[\ce{NO}]^2$.

\textbf{Substitute:}
\[ \text{rate} = k_2K_1[\ce{NO}]^2[\ce{O2}] = k[\ce{NO}]^2[\ce{O2}] \]
where the observed $k$ is the product $k_2K_1$. This matches experiment,
and the intermediate is gone. \checkmark
\end{workedexample}

\begin{apctrap}
The observed rate constant is a \emph{combination} of constants
($k = k_2K_1$ here), which is why a measured $k$ need not correspond to any
single elementary step. That is also why the same rate law can arise from
more than one mechanism.
\end{apctrap}

\segment{INSTRUCTION B}{20}{Multistep energy profiles}

\notesection{Reading a two-hump diagram}{5.10}
\practice{3}

Each elementary step contributes one \answerblank[1.8cm]{hump}. The
number of humps equals the number of \answerblank[2.4cm]{steps}, and the
valley between them is the \answerblank[2.6cm]{intermediate}.

\begin{center}
\begin{tikzpicture}
\begin{axis}[width=10.4cm, height=5cm,
    xlabel={reaction progress}, ylabel={potential energy},
    xtick=\empty, ytick=\empty, xmin=0, xmax=12, ymin=0, ymax=10,
    axis lines=left]
  \addplot[seriesA, samples=250, domain=0:12]
    {2 + 6.5*exp(-((x-3)^2)/1.4) + 4*exp(-((x-8)^2)/1.4)};
  \node[font=\scriptsize,anchor=north] at (axis cs:0.9,1.9) {reactants};
  \node[font=\scriptsize,anchor=south] at (axis cs:3,8.8) {TS 1};
  \node[font=\scriptsize,anchor=north] at (axis cs:5.5,2.4)
    {intermediate};
  \node[font=\scriptsize,anchor=south] at (axis cs:8,6.4) {TS 2};
  \node[font=\scriptsize,anchor=north] at (axis cs:11.2,1.9) {products};
\end{axis}
\end{tikzpicture}
\end{center}

\begin{center}
\fbox{\parbox{0.86\linewidth}{\centering
The \textbf{rate-determining step} is the one with the
\textbf{highest hump} --- the largest activation energy measured from the
species preceding it.}}
\end{center}

In the diagram above, the \answerblank[1.6cm]{first} step is
rate-determining.

\segment{APPLICATION}{20}{Putting the pieces together}

\begin{enumerate}[leftmargin=*,itemsep=0.4em]
  \item For the mechanism $\ce{A + B <=> C}$ (fast equilibrium),
        $\ce{C + B -> D}$ (slow), derive the rate law.
        \workspace[2.6cm]
        \keyonly{\begin{apcanswer}Slow step: rate $= k_2[\ce{C}][\ce{B}]$.
        From step 1, $K_1 = [\ce{C}]/([\ce{A}][\ce{B}])$ so
        $[\ce{C}] = K_1[\ce{A}][\ce{B}]$. Substituting:
        rate $= k_2K_1[\ce{A}][\ce{B}]^2 =
        \mathbf{k[\ce{A}][\ce{B}]^2}$\end{apcanswer}}
  \item A two-step profile shows a first hump of \SI{95}{\kilo\joule} and a
        second of \SI{60}{\kilo\joule} above the intermediate. Which step
        is rate-determining, and what would the rate law depend on?
        \answerlines[3]{The first, since it has the larger barrier. The rate
        law would then be that of the first elementary step, involving only
        the species that collide in it.}
  \item Explain why a valid mechanism can have a rate law containing a
        species that does not appear in the slow step.
        \answerlines[3]{When the slow step involves an intermediate, that
        intermediate is replaced using the equilibrium of an earlier step.
        The substitution brings in the concentrations of the species that
        formed it --- which is how \ce{NO} enters squared in worked
        example 6.}
\end{enumerate}

\begin{exitticket}
Why must an intermediate be eliminated from a rate law, and how is it done?
\keyonly{\begin{apcanswer}Because a rate law must be expressed in
quantities that can actually be measured and controlled, and an
intermediate's concentration is neither --- it is never a species you mix.
It is eliminated using the equilibrium expression of the fast step that
produced it.\end{apcanswer}}
\end{exitticket}

% =====================================================================
\blocklesson{Catalysis}{CED 5.11 \textbullet\ Zumdahl \S12.7}

\begin{objectives}
  \item Explain how a catalyst increases rate.
  \item State precisely what a catalyst does and does not change.
\end{objectives}

\reading{Zumdahl \S12.7}{626--632}

\begin{warmup}[3]
  \item A catalyst is \answerblank[5.6cm]{consumed then regenerated}
  \item The rate-determining step is the one with the:
        \answerblank[3.4cm]{highest barrier}
  \item $k = Ae^{-E_a/RT}$; lowering $E_a$ does what to $k$?
        \answerblank[2.6cm]{increases it}
\end{warmup}

\segment{INSTRUCTION A}{25}{A different route, not a smaller hill}

\notesection{How catalysis works}{5.11}
\practice{6}

A \term{catalyst} provides an \textbf{alternative reaction pathway} with a
lower activation energy. It does not push the reactants over the original
barrier --- it gives them a different, lower one.

\begin{center}
\begin{tikzpicture}
\begin{axis}[width=10.4cm, height=5cm,
    xlabel={reaction progress}, ylabel={potential energy},
    xtick=\empty, ytick=\empty, xmin=0, xmax=10, ymin=0, ymax=10,
    axis lines=left, legend style={draw=none, font=\sffamily\scriptsize,
      at={(0.98,0.98)}, anchor=north east}]
  \addplot[seriesA, samples=200, domain=0:10] {2 + 6.5*exp(-((x-5)^2)/2.0)};
  \addlegendentry{uncatalysed}
  \addplot[seriesB, samples=200, domain=0:10] {2 + 3.2*exp(-((x-5)^2)/2.0)};
  \addlegendentry{catalysed}
  \draw[densely dashed] (axis cs:0,2) -- (axis cs:10,2);
  \node[font=\scriptsize,anchor=west] at (axis cs:0.2,2.5) {reactants};
\end{axis}
\end{tikzpicture}
\end{center}

Notice what the two curves share: the same
\answerblank[2.4cm]{start} and the same \answerblank[2.4cm]{finish}.

\begin{center}
\begin{tabular}{@{}ll@{}}
\hline
\textbf{A catalyst changes} & $E_a$, the mechanism, and therefore the
  \textbf{rate} \\
\textbf{A catalyst does not change} & $\Delta H$, $\Delta G^\circ$, $K$, or
  the position of equilibrium \\
\hline
\end{tabular}
\end{center}

\begin{apctrap}
\textbf{``A catalyst shifts the equilibrium'' is always wrong.} It lowers
$E_a$ for the forward \emph{and} reverse reactions by the same amount, so
both rates increase equally and the equilibrium position is untouched. A
catalyst gets you to the same equilibrium \emph{sooner}; it never moves it.

Equally wrong: ``a catalyst makes the reaction more thermodynamically
favorable.'' $\Delta G^\circ$ is fixed by the reactants and products alone.
\end{apctrap}

\segment{GUIDED PRACTICE}{15}{Types and identification}

\begin{enumerate}[leftmargin=*,itemsep=0.35em]
  \item A \term{homogeneous} catalyst is in the
        \answerblank[3.4cm]{same phase} as the reactants; a
        \term{heterogeneous} catalyst is in a
        \answerblank[3.2cm]{different phase}.
  \item How do you spot a catalyst in a mechanism?
        \answerblank[9.7cm]{it appears in an early step and reappears later}
  \item How do you spot an intermediate?
        \answerblank[7.3cm]{it is produced first, then consumed}
  \item Enzymes are catalysts that are highly
        \answerblank[2.6cm]{specific} to their substrate.
\end{enumerate}

\segment{APPLICATION}{20}{Catalysis reasoning}

\begin{enumerate}[leftmargin=*,itemsep=0.4em]
  \item A catalyst lowers $E_a$ from \SI{120}{\kilo\joule\per\mole} to
        \SI{60}{\kilo\joule\per\mole}. What happens to $\Delta H$, to $K$,
        and to the time needed to reach equilibrium?
        \answerlines[3]{$\Delta H$ and $K$ are both unchanged --- they
        depend only on the reactants and products. The time to reach
        equilibrium falls substantially, because both forward and reverse
        rates increase.}
  \item In the mechanism $\ce{H2O2 + I- -> H2O + IO-}$, then
        $\ce{H2O2 + IO- -> H2O + O2 + I-}$, identify the catalyst and the
        intermediate, and give the overall equation.
        \answerlines[3]{\ce{I-} is the catalyst --- consumed in step 1 and
        regenerated in step 2. \ce{IO-} is the intermediate. Overall:
        $\ce{2H2O2 -> 2H2O + O2}$.}
  \item A student says ``the catalyst was used up, so it must have been a
        reactant.'' What observation would settle it?
        \answerlines[2]{Check whether the species is present again at the
        end. A catalyst is regenerated, so its amount is unchanged overall;
        a reactant is genuinely consumed.}
\end{enumerate}

\begin{exitticket}
Name three things a catalyst changes and three it does not.
\keyonly{\begin{apcanswer}Changes: the activation energy, the mechanism (it
provides a different pathway), and the rate --- forward and reverse alike.
Does not change: $\Delta H$, $K$ and the equilibrium position, or
$\Delta G^\circ$.\end{apcanswer}}
\end{exitticket}

\end{document}
